% Autor: Elias Welt
\section{Datenschicht und API-Integration}
Die gesamte Kommunikation mit dem Backend wird über das Modul \texttt{src/services/api.js} abgewickelt. Diese Methode, oft als "Service Layer Pattern" bezeichnet, hat den Vorteil, dass Komponenten keine Details über URLs oder HTTP-Header kennen müssen.

\subsection{API-Wrapper}
Das \texttt{api}-Objekt fungiert als Wrapper um die native \texttt{fetch}-API. Es implementiert globale Verhaltensweisen:

\begin{enumerate}
    \item \textbf{Basis-URL}: Alle Anfragen sind relativ zu einer zentral konfigurierten Basis-URL
    \item \textbf{Credentials}: Durch \texttt{credentials: 'include'} sendet der Browser automatisch Cookies (für die HttpOnly Session-Cookies) mit jeder Anfrage mit.
    \item \textbf{Content-Type Handling}: JSON-Payloads werden automatisch serialisiert und der korrekte Header gesetzt.
    \item \textbf{Zentrale Fehlerbehandlung}: Antworten mit einem HTTP-Statuscode außerhalb des 200er-Bereichs lösen automatisch eine Exception aus. Dies vereinfacht die Fehlerbehandlung in den Komponenten, da \texttt{try/catch}-Blöcke verwendet werden können.
    \item \textbf{Spezifische Methoden}: Methoden wie \texttt{createChatSession} oder \texttt{uploadDocument} bilden die Logik ab. Letzteres nutzt \texttt{FormData}, um Datei-Uploads (multipart/form-data) zu ermöglichen.
\end{enumerate}

\begin{lstlisting}[language=JavaScript, caption={GET-Methode im API-Service}]
get: async (endpoint, options = {}) => {
  const response = await fetch(`${API_BASE_URL}${endpoint}`, {
    credentials: 'include',
    ...options,
  });
  if (!response.ok) {
    // Fehlerbehandlung: Logging und Werfen eines Errors
    throw new Error(`API Error: ${response.status}`);
  }
  return await response.json();
},
\end{lstlisting}
